\documentclass{article}
\date{}

\usepackage[norsk]{babel}
\usepackage{subcaption}
\usepackage{amsmath}
\usepackage{graphicx}
\usepackage[a4paper, top=1in, bottom=1in, left=1in, right=1in]{geometry}
\usepackage{subcaption}
\usepackage{amsfonts}
\usepackage{dirtytalk}
\usepackage{enumerate}
\usepackage{enumitem}
\usepackage{ulem}

% vertical vector = \vect{...}
\newcommand{\vect}[1]{\begin{pmatrix} #1 \end{pmatrix}}
% double underline is \uu{...}
\newcommand{\uu}[1]{\underline{\underline{#1}}}
% double bold underline is \uub{...}
\newcommand{\uub}[1]{\underline{\underline{\textbf{#1}}}}
% align counter reset
\AtBeginEnvironment{align}{\setcounter{equation}{0}}
% arrows
\newcommand{\lra}{\quad\Longrightarrow\quad}
% equivalent = \eq 
\newcommand{\eq}{\quad\Longleftrightarrow\quad}
\newcommand{\eqq}{\quad=\quad}
% 3x space
\newcommand{\trip}{\ \ \ }
% relational equal
\newcommand{\eqto}[1]{\overset{#1}{\sim}}
% relational equal som støtter linjeskift
\newcommand{\eqqto}[1]{\overset{\substack{#1}}{\sim}}

% linje for å skille oppgaver
\newcommand{\lrule}[1]{%
    \vspace{#1pt} % Space above the line
    \hrule % Draw the horizontal line
    \vspace{#1pt} % Space below the line
}

% matrise med strek, e.g. \begin{pmatrix}[ccc|c]
\makeatletter
\renewcommand*\env@matrix[1][*\c@MaxMatrixCols c]{%
  \hskip -\arraycolsep
  \let\@ifnextchar\new@ifnextchar
  \array{#1}}
\makeatother

\setlength{\parskip}{5pt} 
\setlist{itemsep=8pt, parsep=5pt}


\begin{document}

\subsection*{Oppgave 1}
\(y = \vect{-1\\2\\0}\ \text{ og } \ z=\vect{2\\1\\5} \)
\\\\\\
Løs ligningen \(4y-2x=3z\) for x
\\\\\\
\(2x = 4y-3z \eq \vect{2x_1\\2x_2\\2x_3} = \vect{-4\\8\\0} - \vect{6\\3\\15} \eq \vect{2x_1\\2x_2} = \vect{-10\\5\\15} \eq \vect{x_1\\x_2\\x_3} = \vect{-5\\ 5/2\\ 15/2}\)
\\\\\\
dvs. \(\mathbf{x = \vect{-5\\5/2\\15/2} = d}\)
\\\\\\


\subsection*{Oppgave 2}
\(a = \vect{4, 3, 2}\) og \(b = \vect{-1, -4, -1}\)
\\\\
Prosjekjonen av a ned på b betegnes: \quad \(\mathbf{proj_b A = \frac{a\cdot b}{b \cdot b}b}\)
\\\\
\(a\cdot b = 4\cdot-1 + 3\cdot-4 + 2\cdot-1 = -4-12-2=-18\)
\\\\
\(b\cdot b = -1\cdot-1 + -4\cdot-4 + (-1)\cdot(-1) = 1 + 16 + 1 = 18\)
\\\\
\(proj_bA = -\frac{18}{18}\cdot\vect{-1, -4, -1} = -1 \cdot (-1, -4, -1) = (1, 4, 1)\)
\\\\\\
\(\mathbf{proj_bA = (1, 4, 1) = c}\)
\\\\\\


\subsection*{Oppgave 3}
\(a = \vect{1, 0, 2}\) og \(b = \vect{2, 2, 1}\)
\\\\
\(a \times b = \begin{vmatrix}
    1&1&1\\
    1&0&2\\
    2&2&1
\end{vmatrix} = \begin{vmatrix}
    0&2\\2&1
\end{vmatrix} - \begin{vmatrix}
    1&2\\2&1
\end{vmatrix} + \begin{vmatrix}
    1&0\\2&2
\end{vmatrix} = [(0-4), -(1-4), (2-0)] = [-4, 3, 2]\)
\\\\\\
\(\mathbf{a\times b = (-4, 3, 2) = a}\)
\\\\\

\newpage
\subsection*{Oppgave 4}
\(A = \vect{2&3&-1\\0&2&4}\) og \(x = \vect{1\\1\\2}\)
\\\\
\(Ax = \vect{2&3&-1 \\ 0&2&4}\vect{1\\1\\2} = \begin{pmatrix}
    2+3-2 \\ 0+2+8
\end{pmatrix} = \begin{pmatrix}
    3 \\ 10
\end{pmatrix}\)
\\\\\\
\(\mathbf{Ax = \vect{3 \\ 10} = b}\)
\\\\\\

\subsection*{Oppgave 5}
\(A = \vect{1&2&4\\0&2&3\\0&0&3}\) og \(B = \vect{2&0&0\\1&4&0\\2&7&4}\)
\\\\\\
\(AB = \begin{pmatrix}
    2+2+8 & 0+8+28 & 0+0+16 \\
    0+2+6 & 0+8+21 & 0+0+12 \\
    0+0+6 & 0+0+21 & 0+0+12
\end{pmatrix} = \begin{pmatrix}
    12 & 36 & 16 \\
    8 & 29 & 12 \\
    6 & 21 & 12
\end{pmatrix}\)
\\\\\\
\(\mathbf{AB = \begin{pmatrix}
    12 & 36 & 16 \\
    8 & 29 & 12 \\
    6 & 21 & 12
\end{pmatrix} = b}\)
\newpage

\subsection*{Oppgave 6}
\(A, B\ \in \mathbb{R}^{n\times n}\)
\\\\\\
\((B^TA)^T(AB)^{-1}\)
\\\\
\((B^TA)^T = A^TB\)  \quad og \quad \((AB)^{-1} = (B^{-1}A^{-1})\)
\\\\
\(A^TB(B^{-1}A^{-1}) = A^T((BB^{-1})A^{-1}) = A^T(I_nA^{-1}) = A^T(A^{-1}) = A^TA^{-1}\)
\\\\\\
\(\mathbf{(B^T)A^T(AB)^{-1} = A^TA^{-1} = e}\)
\\\\\\


\subsection*{Oppgave 7}
\(\begin{cases}
    2x+y+3z=0\\
    x+4z=a\\
    5x+2y+10z=2
\end{cases}\)
\\\\\\
\(\begin{pmatrix}
    2&1&3&0\\
    1&0&4&a\\
    5&2&10&2
\end{pmatrix}
\eqqto{I \leftrightarrow II}
\begin{pmatrix}
    1&0&4&a\\
    2&1&3&0\\
    5&2&10&2
\end{pmatrix}
\eqqto{II-2I\\\\III-5I}
\begin{pmatrix}
    1&0&4&a\\
    0&1&-5&-2a\\
    0&2&-10&2-5a
\end{pmatrix}
\eqqto{III-2II}
\begin{pmatrix}
    1&0&4&a\\
    0&1&-5&-2a\\
    0&0&0&2-a
\end{pmatrix}
\\\\\\
\begin{cases}
    x + 4z = a \eq x = a-4z \eq x = 2-4z\\
    y -5z = -2a \eq y = 5z-2a \eq y = 5z-4\\
    0 = 0 \eq z = z \eq z \in \mathbb{R}
\end{cases}
\\\\\\
\text{ og for \textbf{a = 2 = b)}  finnes det uendelig mange løsninger.}
\)
\\\\
\(\mathbf{a = 2 = b})\)
\\\\


\subsection*{Oppgave 8}
\(
A = \vect{2&2&6 \\ 2&1&4 \\ 1&2&5 \\ 0&1&2}
\eqqto{II-I}
\vect{2&2&6 \\ 0&-1&-2 \\ 1&2&5 \\ 0&1&2}
\eqqto{I/2}
\vect{1&1&3 \\ 0&-1&-2 \\ 1&2&5 \\ 0&1&2}
\eqqto{III-I}
\vect{1&1&3 \\ 0&-1&-2 \\ 0&1&2 \\ 0&1&2}
\eqqto{IV-III\\\\I+II}
\vect{1&0&1 \\ 0&-1&-2 \\ 0&1&2 \\ 0&0&0}
\\\\\\
\eqqto{III+II}
\vect{1&0&1 \\ 0&-1&-2 \\ 0&0&0 \\ 0&0&0}
\eqqto{-1II}
\vect{1&0&1 \\ 0&1&2 \\ 0&0&0 \\ 0&0&0}
\mathbf{ = a)}\)
\\\\\\


\newpage
\subsection*{Oppgave 9}
\(\vect{-2&-6&-3&2 \\ 1&3&2&-1 \\ 2&6&4&0}
\eqqto{I+II}
\vect{-1&-3&-1&1 \\ 1&3&2&-1 \\ 2&6&4&0}
\eqqto{II+I}
\vect{-1&-3&-1&1 \\ 0&0&1&0 \\ 2&6&4&0}
\eqqto{III+2I}
\vect{-1&-3&-1&1 \\ 0&0&1&0 \\ 0&0&-2&-2}
\eqqto{III+2II}
\\\\
\vect{-1&-3&-1&1 \\ 0&0&1&0 \\ 0&0&0&-2}
\eqqto{-1I\\\\-III/2}
\vect{1&3&1&-1 \\ 0&0&1&0 \\ 0&0&0&1}
\eqqto{I-II\\\\I+III}
\vect{1&3&0&0 \\ 0&0&1&0 \\ 0&0&0&1}
\)
\\\\\\
Søylekolonner: \qquad \textbf{kolonne 1, 3 og 4 = d)}
\\


\subsection*{Oppgave 10}
\(\vect{1&2&4 \\ 2&1&3 \\ 2&7&13} x = b\)
\\\\\\
\textbf{hint:} \\\\
\(
\begin{pmatrix}[ccc|c]
1&2&4&b_1 \\ 
2&1&3&b_2 \\ 
2&7&13&b_3
\end{pmatrix}
\eqqto{II-2I\\\\III-2I}
\begin{pmatrix}[ccc|c]
1&2&4&b_1 \\ 
0&-3&-5&b_2-2b_1 \\ 
0&3&5&b_3-2b_1
\end{pmatrix}
\eqqto{III+II}
\begin{pmatrix}[ccc|c]
1&2&4&b_1 \\ 
0&-3&-5&b_2-2b_1 \\ 
0&0&0&b_3+b_2-4b_1
\end{pmatrix}
\eqqto{-II/3}
\begin{pmatrix}[ccc|c]
1&2&4&b_1 \\ 
0&1&5/3&-(b_2-2b_1)/3 \\ 
0&0&0&b_3+b_2-4b_1
\end{pmatrix}
\)
\\\\\\
For at en løsning skal finnes så må $(b_3+b_2-4b_1) = 0$ 
\\\\\\
Dette tilfredsstilles hvis $\mathbf{b = \vect{1\\3\\1} = e)}$
\\\\\\
Da blir $(b_3+b_2-4b_1) = (1+3-4) = 0$

\newpage
\subsection*{Oppgave 11}
Hvilken påstand er sann?
\begin{enumerate}
    \item Det finnes systemer $Ax=b$ som har nøyaktig to løsninger.

    \item Hvis A har et pivotelement i hver kolonne så har $Ax=b$ alltid løsning(er)\\

    \item Trappeformen til en matrise er unik

    \item $Ax=b$ har alltid løsning(er) når A har flere kolonner enn rader

    \item \textbf{Hvis A har et pivotelement i hver rad så har $Ax = b$ alltid løsninger}\\
    - Dette er sant fordi vet at hver 
    
\end{enumerate} 


\subsection*{Oppgave 12}
\(A =\begin{cases}
    2x_1+x2=1-3x_3\\
    x_2+4x_3=x_1-12\\
    3x_1+x_3=5
\end{cases}
=
\begin{cases}
    2x_1+x_2+3x_3=1\\
    -x_1+x_2+4x_3=-12\\
    3x_1+x_3=5
\end{cases}\)
\\\\\\
\(A' = 
\begin{pmatrix}[ccc|c]
    2&1&3&1\\
    -1&1&4&-12\\
    3&0&1&5
\end{pmatrix}
\eqqto{I \leftrightarrow II\\\\-1I}
\begin{pmatrix}[ccc|c]
    1&-1&-4&12\\
    2&1&3&1\\
    3&0&1&5
\end{pmatrix}
\eqqto{II-2I\\\\III-3I}
\begin{pmatrix}[ccc|c]
    1&-1&-4&12\\
    0&3&11&-23\\
    0&3&13&-31
\end{pmatrix}
\eqqto{III-II}
\begin{pmatrix}[ccc|c]
    1&-1&-4&12\\
    0&3&11&-23\\
    0&0&2&-8
\end{pmatrix}
\\\\\\
\eqqto{(1/2)III}
\begin{pmatrix}[ccc|c]
    1&-1&-4&12\\
    0&3&11&-23\\
    0&0&1&-4
\end{pmatrix}
\eqqto{(1/3)II}
\begin{pmatrix}[ccc|c]
    1&-1&-4&12\\
    0&1&11/3&-23/3\\
    0&0&1&-4
\end{pmatrix}
\eqqto{I+II\\\\I+(1/3)III\\\\II-(11/3)III}
\begin{pmatrix}[ccc|c]
    1&0&0&3\\
    0&1&0&7\\
    0&0&1&-4
\end{pmatrix} \\
x = \vect{3\\7\\-4} = e
\)
\\\\\\\\

\subsection*{Oppgave 13}
\(T:\mathbb{R}³ \rightarrow \mathbb{R}²\) er lineær avbildning slik at 
\\\\
\(T(e_1)=\vect{1\\2}, \qquad T(e_2)=\vect{3\\-3}, \qquad T(e_3)=\vect{0\\5}\)
\\\\
Der $e_k$ (for k = 1,2,3) er i vektoren $\mathbb{R}³$ med k-te komponent lik 1, og med alle andre komponenter lik 0.
\\\\
Hva kan vi si om T?

\begin{enumerate}[a)]

    \item $T(x)=\vect{1&2\\3&-3\\0&5} x$ for alle $x\in\mathbb{R}^3$
    
    \item T er inverterbar, og $T^{-1} \vect{\vect{3\\-3}} = e_2$
    
    \item $\T(x)=\vect{1&3&0\\2&-3&5}x$ for alle $x\in\mathbb{R}^3$\}\\
    - Riktig
    
    \item $T=\vect{1&0&0\\0&1&0\\0&0&0}$


    \item $T(x)\cdot T(y) = T(x\cdot Y)$ for alle $x,y \in \mathbb{R}^3$
    
\end{enumerate}


\subsection*{Oppgave 14}
\(A \in \mathbb{R}^2\) er mengden av alle punkter i trekanten som har hjørner i punktene $(1, 1), (2, 1), (2, 2)$
\\\\
Hvilket utsagn er sant?

\([1,2]\times[1,2], x\geq y \eqq \{(1, 1), (2, 1), (2,2)\} \eqq a\)


\subsection*{Oppgave 15}

\subsection*{Oppgave 16}

\(
\forall p, P(p), \quad \exists q, P(q), \quad q > p 
\\\\\\
\neg(\forall p, P(p), \quad \exists q, P(q), \quad q > p)
\\\\
\exists p, P(p), \quad \neg(\exists q, P(q), \quad q > p)
\\\\
\exists p, P(p), \quad \forall q, P(q), \quad \neg(q > p)
\\\\
\exists p, P(p), \quad \forall q, P(q), \quad q \leq p
\)
\\\\\\
Dvs. \quad Det finnes et primtall p for alle primtall q hvor q $\leq$ p
\\\\
Det finnes et primtall p slik at $q \leq p$ q for alle primtall q. = d

\subsection*{Oppgave 17}

"Det finnes et reelt tall a slik at x+a = x for alle relle tall x."
\\\\
\(\)

\end{document}
